\section*{Supplementary Material}
For both experiments, toxicity estimation of PFAS molecules \cite{feinstein2021uncertainty} \footnote{AI4PFAS project \url{ https://github.com/AI4PFAS/AI4PFAS.}} and biodegradability of general molecules \cite{lee2022comparative}, we divided the datasets into: 60\% for training purposes, 20\% for validation purposes, and 20\% for testing purposes. In both cases, the divisions have been made following the original papers in order to have a fair comparison with the state-of-the-art algorithms.
Table \ref{TableTox} refers to the hyper-parameters considered for the toxicity estimation task. 

\begin{table}[htbp]
\scriptsize
\centering
\caption{\textsc{MultiModal}-\textsc{MoLFormer} hyper-parameters for toxicity prediction considering the LDToxDB dataset.}

\label{TableTox}
\begin{tabular}{@{}p{4cm}p{2cm}}
\toprule
\textbf{Hyper-parameter} & \textbf{Value}  \\ \midrule

Device & cuda  \\
GPU & 1 \\
Batch Size & 32  \\
Number of heads & 12  \\
Number of embeddings & 768  \\
Dropout & 0.1  \\
Learning rate & 3e-5\\
Number of workers & 8\\
Number of epochs & 2000\\
Dimension of the Feed-Forward network & [797 797 797 1] \\
 
 \\\bottomrule
\end{tabular}
\end{table}
%%%%%%%%%%%%%%%%%%%%%%%%%%%%%%%%%%%%%%%%%%%%%%%%%%%%%%%%%%%%

Table \ref{TableBiodeg} refers to the hyper-parameters considered for biodegradability classification.. 


\begin{table}[htbp]
\scriptsize
\centering
\caption{\textsc{MultiModal}-\textsc{MoLFormer} hyper-parameters for biodegradability prediction.}

\label{TableBiodeg}
\begin{tabular}{@{}p{4cm}p{2cm}}
\toprule
\textbf{Hyper-parameter} & \textbf{Value}  \\ \midrule

Device & cuda  \\
GPU & 1 \\
Batch Size & 32  \\
Number of heads & 12  \\
Number of layers & 12  \\
Number of embeddings & 768  \\
Dropout & 0.1 \\
Learning rate & 3e-5\\
Number of workers & 8\\
Number of epochs & 2000\\
Dimension of the Feed-Forward network & [773 773 773 1]\\

\\\bottomrule
\end{tabular}
\end{table}

Both experiments have been conducted in the IBM computing cluster with $32$ GB of RAM memory and a  Nvidia V$100$ GPU.

Fig. \ref{fig:evol} illustrates the evolution of the \textsc{MultiModal}-\textsc{MoLFormer} in terms of classification accuracy for the toxicity estimation task considering different numbers of training samples. Through this analysis, we could observe that the proposed \textsc{MultiModal}-\textsc{MoLFormer} needed $4096$ samples to obtain similar classification accuracy to the original \textsc{MoLFormer} approach that used the full training set. This indicate us, the influence of the physico-chemical features to improve the learning of the proposed approach, even considering less training data samples. 

\begin{figure}[htbp]
    \centering
    \includegraphics[width=12cm]{supp2.png}
    \caption{
        This figure illustrates the \textsc{MultiModal}-\textsc{MoLFormer} learning evolution with different training data samples sizes. 
    }
    \label{fig:evol}
\end{figure}

The experiments conducted in this research are based on the recently introduced \textsc{MoLFormer} approach which can be accessed through this link (\url{https://github.com/IBM/molformer}). 

To extract the physico-chemical features based on the Mordred descriptors we used the library available at \url{http://mordred-descriptor.github.io/documentation/v0.1.0/introduction.html} \cite{moriwaki2018mordred}. To build the Markov-Blanket causal feature selection blocks we used the library available at \url{https://github.com/atif-hassan/PyImpetus} \cite{hassan2021ppfs}. The default parameters were used for the physico-chemical extraction from SMILES and to build the Markov-Blanket causal graphs.



